\section{Love and Strife}

\begin{quotex}
Friends, I know indeed that truth is in the word I shall utter; but it is hard for men, and jealous are they of the assault of belief in their souls.
\flright{\textsc{Empedocles}}
\end{quotex}

It was difficult to put aside the bottle of El Padrino anejo\footnote{\url{http://tequila1921.com/web/en/products.html}} and a beckoning Cohiba to write this tonight, given the level of meaningless strife recently based on near total incomprehension and ignorance. But strife is not to be avoided, but rather overcome, first in oneself. So I am curious about what in Gornahoor is causing confusion. We have made clear that paganism arose without reference to or concern about Christianity, unlike neo-paganism which is a contemporary artificially constructed system based on a romanticized version of Odinism, a crude grasp of Nietzsche, and a rabid anti-Chrisitianism based more on Enlightenment criticism than any Traditional perspective.

Since there is less known of the Northern tradition, especially in any metaphysical way, we have been emphasizing the better documented Greek and Roman traditions. Guenon simply assumed the Greeks thinkers were inferior to the Asians, so he neglected it and Evola simply assumed that any properly educated man would know the historical background, so he never fully explicated it. Both assumptions are false, though in different ways.

There is real wisdom and an authentic tradition in paganism (properly understood) which we have endeavored to bring out, given the limitations of a blog. Alongside the rites of the city, there were the cults of the Olympian gods and the various mystery religions. Later, the philosophers would attempt in various ways to explicate the ancient wisdom verbally, signaling the shift from mythos to logos. One of the more interesting systems was that of Empedocles who builds on \textbf{Orpheus},
\textbf{Pythagoras}, and \textbf{Parmenides}.

We know of Empedocles from large fragments of two poems which Aristotle described as esoteric. Man in the primordial state was free of Strife under the reign of Aphrodite, the Queen of Love. God is Love as well as Logos and Law. Yet man fell from this state of innocence because of a primal sin, which created the world of strife. The sin was animal sacrifice, a rite of the ancestor worship. This is the path of darkness, or pitriyana, counter to which Empedocles offers the path of light, or devayana.

It begins with a stage of purification. Then, theoria, when the seeker learns the doctrine in an intellectual way, trying to understand it as a whole from the various parts. Finally, one achieves theosis, the state of the god-man; Empedocles himself claimed to be divine. This path is the same as in the mystery religions and is still followed today by the monks on Mt. Athos. This is summarized in `Blessed are the pure in heart, for they shall see God', since, according to Empedocles, only the like can know like. That some see in this doctrine merely a concession to weakness is just another sign of the times.

Empedocles teaches the cycle of existence\footnote{\url{https://gornahoor.net/?p=1863}}. In our cycle, we pass from the \textbf{Reign of Love} which is unity and harmony to the \textbf{Reign of Strife}, which is separation and evil. This cycle ends when Love is freed from Strife and brought back to the unity of God.

For further reference:

\emph{Reality} by Peter Kingsley

\emph{From Religion to Philosophy} by F. M. Cornford

As for the Mt Athos reference, a manuscript of missing parts of Empedocles' poems was discovered in a monastery in 1851.


\flrightit{Posted on 2011-09-07 by Cologero}

\begin{center}* * *\end{center}

\begin{footnotesize}
\begin{sffamily}
\texttt{Graham on 2011-09-07 at 09:22 said:}

`Such is the life of the gods reign of love ; but of other souls, that which follows God best and is likest to him like knows like, knowledge is being lifts the head of the charioteer into the outer world, and is carried round in the the revolution, troubled indeed by the steeds, and with difficulty beholding true being; while another rises and falls, and sees, and again fails to see by reason of the unruliness of the steeds. The rest of the souls are also longing after the upper world and they all follow, but not being strong enough they are carried round below the surface, plunging, treading on one another, each striving to be first reign of strife ; and there is confusion and perspiration and extremity of effort; and many of them are lamed or have their wings broken through the ill-driving of the charioteers; and all of them after a fruitless toil, not having attained to the mysteries of true being, go away, and feed upon opinion.'


\hfill

\end{sffamily}
\end{footnotesize}

